\section{Polytopes}

A polytope is the general name for a regular shape in any dimension, a polygon in two, a solid in three, and onward. Four dimensions has exactly six regular ones, the 5-cell, the 8-cell or tesseract, the 16-cell, the 24-cell, the 120-cell, and the 600-cell. Five of them are higher-dimensional versions of three-dimensional solids, the way the tesseract is the four-dimensional cube. The 24-cell is not a version of anything. It is self-dual, meaning it has the same number of corners as cells, twenty-four of each, with every cell an octahedron, and it has no counterpart in three dimensions or in any other. Table~\ref{tab:polytopes} lays out the ladder from the segment up to five dimensions.

\begin{table}[ht]
\centering
\small
\begin{tabular}{@{}rll@{}}
\toprule
dimension & the figure & the regular ones \\
\midrule
$1$ & segment & the line segment \\
$2$ & polygon & triangle, square, pentagon, and on \\
$3$ & polyhedron & tetrahedron, cube, octahedron, dodecahedron, icosahedron \\
$4$ & polychoron & 5-cell, tesseract, 16-cell, 24-cell, 120-cell, 600-cell \\
$5$ & 5-polytope & simplex, hypercube, cross-polytope \\
\bottomrule
\end{tabular}
\caption{The regular polytopes by dimension, from the segment to five. The figure's general name changes at each step, and the 24-cell, in four dimensions, is the one with no analog above or below.}
\label{tab:polytopes}
\end{table}

Place the 24-cell's corners at the natural coordinates and they become a famous set. They are the vectors with two entries $\pm 1$ and two entries zero, twenty-four in all, the root system $D_4$, a highly symmetric set of vectors that sits at the heart of a major family of symmetries. They are also the twenty-four unit Hurwitz quaternions, the smallest such set closed under multiplication, where quaternions are four-dimensional numbers that encode rotations, and this particular set forms a group called $2T$. So the twenty-four directions out of a dock are at once a symmetric set of vectors and a group of rotations, and to compose two of them is just to multiply two quaternions, exactly and with no rounding.

That last fact carries the deepest gift of the structure. Quaternion multiplication knows the difference between a half turn and a full one.

\begin{theorem}[Double cover]
On the twenty-four directions a rotation by $2\pi$ is represented by $-1$, and a rotation by $4\pi$ by the identity. This is the defining mark of a spinor, and it holds exactly in the finite group $2T$, with no limit taken.
\end{theorem}

A spinor is a thing that must be turned around twice to come back to itself, and the electron is one. Here that behavior is not added by hand or reached for in a continuum. It is already present in the multiplication table of twenty-four directions, a fact about a finite group of order twenty-four.

The twenty-four directions hold one more structure. Under the full symmetry they split into three octets, written $8v$, $8s$, and $8c$, a vector octet and two spinor octets of opposite handedness. The three are interchanged by an exact threefold symmetry called triality, which exists for this structure and no other. Each octet will earn a physical role far downstream, the vector becoming light and the two spinors the two handednesses of matter, with the threefold symmetry becoming the three generations. Here it is only a fact about the 24-cell.

\begin{proposition}[Triality]
The twenty-four roots of $D_4$ split into three octets $8v$, $8s$, $8c$, permuted by an order-three triality symmetry. Lifting triality to its natural algebra gives the exceptional Jordan algebra $J_3(\mathbb{O})$, whose rank is forced to three by the non-associativity of the octonions.
\end{proposition}

The 24-cell also sits on a famous climb. Its two larger siblings, the 120-cell and the 600-cell, are built on the golden ratio and carry the symmetry of the icosahedron into four dimensions. The 24-cell stands one rung below them, simpler and self-dual, and one rung above the cube. Stack 24-cells the right way and their directions extend through the exceptional groups in order, $D_4$ to $F_4$ to $E_6$, $E_7$, and finally $E_8$, the densest and most symmetric lattice there is. The base sits at the foot of that climb, on the smallest rung that already carries spin, which is exactly as high as it needs to go and no higher.
